\begin{abstract}
Looped language models apply a shared stack of layers several times per token, so a single
model can trade inference compute for accuracy by changing how many loops it runs. This makes
them a natural fit for agents, whose requests range from single tool calls to long chains of
dependent calls. In practice, however, deployed looped models run a fixed number of loops for
every request, and per-token exit gates lose most of their savings once requests are decoded in
batches. Can a looped model decide, request by request, when extra loops are worth their cost?
We study depth allocation in looped models for multi-step tool use and propose an escalation
router: every request starts at a low loop count, and the model switches to full depth only when
its own early loops disagree on the decision tokens of its first tool call, a signal it computes
anyway. A controlled analysis motivates this design: choosing depth separately for each tool call
gives no gain beyond a sampling control with the same budget, whereas the depth a request needs
is a stable property that early loops reveal and the request text does not. On held-out NESTful
requests with a fine-tuned Ouro-2.6B, the router \pending{matches full-depth accuracy at X\% of
its compute}. We also show that most MathQA validation problems have near-duplicates in its
training split, a source of inflated results for work that trains on MathQA, and provide clean
splits.
\end{abstract}
